\chapter*{Resumo}
%\addcontentsline{toc}{chapter}{Resumo}

A fala sussurrada, caracterizada pela falta de vibração das pregas vocais, é frequentemente utilizada para comunicação discreta, mas apresenta desafios sociais e profissionais significativos para indivíduos com afonia que dependem dela como seu modo principal de fala. Métodos de reabilitação existentes, como o eletrolaringe, a fala esofágica e a tecnologia de conversão de texto em fala, proporcionam algum alívio, mas não oferecem soluções práticas e de som natural. Este documento apresenta uma abordagem inovadora para converter a fala sussurrada em fala vocalizada de som natural. A principal motivação por trás desta pesquisa é melhorar a comunicação para indivíduos que dependem da fala sussurrada devido a condições médicas. O sistema proposto integra a geração de pulsos glotais e a modelagem do trato vocal para realizar a conversão de fala em tempo real, visando melhorar a inteligibilidade e a naturalidade da fala sussurrada.

O sistema foi projetado para operar em tempo real numa plataforma baseada em STM32, demonstrando a viabilidade da conversão de fala sussurrada para vocalizada em tempo real sem atrasos significativos. Componentes-chave incluem algoritmos para a geração de pulsos glotais, que fornecem o sinal de excitação vocal, e a modelagem do trato vocal utilizando {\it Linear Predictive Coding} (LPC). A arquitetura garante uma síntese de fala suave e contínua ao interpolar filtros entre pulsos individuais e utilizando um contorno de F0 para imitar a prosódia natural. Testes perceptivos abrangentes foram realizados para avaliar o desempenho do sistema, incluindo testes de conversão de fala sussurrada para vocalizada e avaliações de síntese de vogais realizadas em um ambiente web MUSHRA. Os resultados mostraram que o sistema completo supera os mecanismos de filtragem de pulsos individuais e de quadros sem segmentação em termos de inteligibilidade e naturalidade da fala. A análise estatística confirmou a significância dessas melhorias, destacando a eficácia da abordagem proposta de síntese híbrida no domínio tempo-frequência.

Além das contribuições técnicas, a tese alinha-se com vários Objetivos de Desenvolvimento Sustentável (ODS) das Nações Unidas, tais como promover a educação inclusiva, o crescimento econômico e a redução das desigualdades. Ao facilitar uma melhor comunicação para indivíduos com deficiências vocais, o sistema apoia a inclusão social e a acessibilidade em diversos contextos.

\chapter*{Abstract}
%\addcontentsline{toc}{chapter}{Abstract}

Whispered speech, characterized by the lack of vocal fold vibration, is often used for discreet communication but poses significant social and professional challenges for individuals with aphonia who rely on it as their primary mode of speech. Existing rehabilitation methods such as the electrolarynx, esophageal speech, and text-to-speech technology provide some relief but fall short of offering practical and natural-sounding solutions. This document presents a novel approach for converting whispered speech into natural-sounding voiced speech. The primary motivation behind this research is to enhance communication for individuals who rely on whispered speech due to medical conditions. The proposed system integrates glottal pulse generation and vocal tract filter modeling to perform on-the-fly speech conversion, aiming to improve the intelligibility and naturalness of whispered speech.

The system is designed to operate in realtime on an STM32-based platform, demonstrating the feasibility of on-the-fly whispered-to-voiced speech conversion without significant delays. Key components include algorithms for glottal pulse generation, which provide the voiced excitation signal, and vocal tract filter modeling using Linear Predictive Coding (LPC) analysis. The architecture ensures smooth and continuous speech synthesis by interpolating filters between individual pulses and employing an F0 contour to mimic natural prosody. Comprehensive perceptual testing was conducted to evaluate the system's performance, including whispered to speech conversion tests and vowel synthesis assessments performed in a MUSHRA web environment. The results showed that the complete system outperforms individual pulse and frame filtering mechanisms without segmentation in terms of speech intelligibility and naturalness. Statistical analysis confirmed the significance of these improvements, highlighting the effectiveness of the proposed hybrid time-frequency domain synthesis approach.

In addition to the technical contributions, the thesis aligns with several United Nations Sustainable Development Goals (SDGs), such as promoting inclusive education, economic growth, and reducing inequalities. By facilitating better communication for individuals with voice impairments, the system supports social inclusion and accessibility in various contexts.


